% =====================================================================
%  Projet 10 — Self-supervised learning contrastif (SimCLR) sur CIFAR-10
%  Cours 14, Projets Master — TCHAYE-KONDI Jude, Ph.D.
%
%  COMPILATION : xelatex rapport_projet10.tex  (deux fois, pour les renvois)
%
% =====================================================================
\documentclass[11pt,a4paper]{article}

% --- noms des membres : un seul endroit a modifier ---------------------
\newcommand{\membreKR}{KAMBIA Rafiatou}
\newcommand{\membreAK}{ALLAGLO Kossiko}
\newcommand{\membreAA}{AMEGNRAN Athanase}
\newcommand{\lienDepot}{https://github.com/sunborndev/simclr-cifar10}

\usepackage{fontspec}
\usepackage{polyglossia}
\setmainlanguage{french}
\setmainfont{TeX Gyre Pagella}
\setsansfont{TeX Gyre Heros}
\setmonofont[Scale=0.86]{DejaVu Sans Mono}

\usepackage[a4paper,top=2.4cm,bottom=2.4cm,left=2.5cm,right=2.5cm]{geometry}
\usepackage{amsmath,amssymb}
\usepackage{graphicx}
\usepackage{booktabs}
\usepackage{tabularx}
\usepackage{multirow}
\usepackage{xcolor}
\usepackage{caption}
\usepackage{enumitem}
\usepackage{titlesec}
\usepackage{fancyhdr}
\usepackage{microtype}
% commande longue : \brk marque les points de coupure autorises
\newcommand{\brk}{\discretionary{}{}{}}
\newcommand{\cmd}[1]{\texttt{\small #1}}
\emergencystretch=3em
\hyphenpenalty=200
\tolerance=1200
\usepackage[colorlinks=true,linkcolor=bleuP,citecolor=bleuP,urlcolor=bleuP]{hyperref}

\definecolor{bleuP}{HTML}{2a78d6}
\definecolor{aquaP}{HTML}{1baf7a}
\definecolor{orangeP}{HTML}{eb6834}
\definecolor{encre2}{HTML}{52514e}
\definecolor{fondgris}{HTML}{f5f5f3}

\captionsetup{font=small,labelfont={bf,color=bleuP},width=0.95\textwidth}
\titleformat{\section}{\large\bfseries\color{bleuP}}{\thesection}{0.7em}{}
\titleformat{\subsection}{\normalsize\bfseries}{\thesubsection}{0.6em}{}
\titlespacing*{\section}{0pt}{1.6em}{0.6em}
\titlespacing*{\subsection}{0pt}{1.1em}{0.35em}
\setlist{itemsep=1pt,topsep=3pt,parsep=0pt}

\pagestyle{fancy}\fancyhf{}
\fancyhead[L]{\footnotesize\color{encre2}Projet 10 — SSL contrastif (SimCLR) sur CIFAR-10}
\fancyhead[R]{\footnotesize\color{encre2}\thepage}
\renewcommand{\headrulewidth}{0.3pt}

% encadre "a retenir"
\newcommand{\encadre}[1]{%
  \par\vspace{4pt}\noindent
  \colorbox{fondgris}{\parbox{\dimexpr\textwidth-2\fboxsep}{\small #1}}%
  \par\vspace{4pt}}

\newcommand{\fig}[4]{%
  \begin{figure}[htbp]\centering
  \includegraphics[width=#2\textwidth]{#1}
  \caption{#3}\label{#4}
  \end{figure}}

\begin{document}

% =====================================================================
\begin{center}
{\LARGE\bfseries Apprentissage contrastif auto-supervisé\\[2pt]
pour le pré-entraînement visuel :}\\[6pt]
{\large de combien d'annotations SimCLR nous dispense-t-il sur CIFAR-10 ?}\\[14pt]

{\large \membreKR{} \quad\textbullet\quad \membreAK{} \quad\textbullet\quad \membreAA{}}\\[4pt]
{\small Master --- Cours 14, Projets de recherche appliquée\\
Encadrant : TCHAYE-KONDI Jude, Ph.D. \quad\textbullet\quad Août 2026}\\[4pt]
{\small Code et instructions de reproduction : \url{\lienDepot}}
\end{center}
\vspace{6pt}

\noindent\colorbox{fondgris}{%
\parbox{\dimexpr\textwidth-2\fboxsep}{\small
\textbf{Résumé.} Annoter des images coûte cher ; le pré-entraînement auto-supervisé promet
d'exploiter les images non étiquetées pour réduire ce coût. Nous implémentons un cadre
contrastif de type SimCLR en PyTorch, pré-entraînons un ResNet-18 sur les 50\,000 images
de CIFAR-10 \emph{sans aucun label}, puis mesurons ce que ce pré-entraînement apporte
lorsque 1\,\%, 10\,\% et 100\,\% des annotations sont disponibles. La comparaison est
appariée : le fine-tuning de l'encodeur pré-entraîné et la baseline supervisée
\emph{from scratch} partagent la totalité de leurs hyperparamètres, seuls les poids
initiaux diffèrent. Avec 450 images étiquetées, le pré-entraînement fait passer
l'accuracy de test de $39{,}52\,\%$ à $72{,}10\,\%$, soit $+32{,}58$ points ; avec
45\,000 labels, les deux méthodes sont à parité. Un bras de contrôle à encodeur aléatoire
gelé ($25{,}66\,\%$) établit que ce gain vient du pré-entraînement et non de
l'architecture. Une étude d'ablation additive sur cinq combinaisons d'augmentations
montre que les perturbations de couleur portent l'essentiel de la qualité de la
représentation ($+22{,}54$ points à elles deux), tandis que le retournement horizontal et
le flou gaussien n'ont aucun effet mesurable à cette échelle. Nous documentons enfin une
dissociation nette entre la qualité de la tâche prétexte et celle de la représentation :
la configuration qui minimise le mieux la perte NT-Xent produit la pire représentation. Un
diagnostic mené sur les données seules, sans réseau de neurones, en donne le mécanisme
chiffré : privé de distorsion de couleur, un simple histogramme de couleur retrouve la
bonne image dans $51{,}9\,\%$ des cas parmi 2\,000 candidates, contre $0{,}05\,\%$ au
hasard --- la tâche prétexte est à moitié résolue avant tout apprentissage.
\medskip

\noindent\textbf{Mots-clés :} apprentissage auto-supervisé, apprentissage contrastif,
SimCLR, NT-Xent, étude d'ablation, régime de faible supervision.
}}

% =====================================================================
\section{Introduction}

\paragraph{Contexte.} Les réseaux de neurones profonds appliqués à la vision supposent
généralement l'existence d'un jeu de données annoté en quantité suffisante. En pratique,
cette hypothèse est la plus coûteuse du pipeline : l'annotation demande du temps humain,
parfois une expertise rare (imagerie médicale, télédétection), et devient impraticable à
l'échelle des flux d'images réellement disponibles. Les images \emph{non étiquetées}, en
revanche, sont abondantes et quasi gratuites.

\paragraph{Problématique.} L'apprentissage auto-supervisé (\emph{self-supervised
learning}, SSL) propose de construire un signal d'apprentissage à partir des données
elles-mêmes, sans intervention humaine. Dans le cadre contrastif de SimCLR
\cite{chen2020simclr}, ce signal est simple : deux vues augmentées de la \emph{même} image
doivent produire des représentations proches, et se distinguer des vues des autres images
du lot. Le modèle apprend ainsi ce qui, dans une image, résiste aux déformations que l'on
lui inflige --- et c'est précisément ce qui définit son contenu sémantique.

La question que nous traitons est quantitative et non conceptuelle :

\encadre{\textbf{Question de recherche.} À budget de calcul strictement égal, de combien
d'images annotées un pré-entraînement contrastif dispense-t-il un classifieur CIFAR-10 ?
Et quelles augmentations de données ce pré-entraînement exige-t-il pour produire une
représentation utile ?}

\paragraph{Contributions.} Ce travail apporte les six éléments suivants.
\begin{enumerate}[label=(\roman*)]
\item Une implémentation de bout en bout du cadre SimCLR en PyTorch --- construction des
paires, encodeur partagé, tête de projection, perte NT-Xent --- écrite à partir des
équations du papier et non reprise d'une bibliothèque de haut niveau (section~\ref{sec:methode}).
\item Une comparaison \emph{appariée} sur trois régimes d'annotation (450, 4\,500,
45\,000 labels) opposant quatre bras : fine-tuning du pré-entraîné, évaluation linéaire du
pré-entraîné gelé, baseline supervisée \emph{from scratch}, et un témoin à encodeur
aléatoire gelé (section~\ref{sec:resultats}).
\item Une étude d'ablation additive des augmentations en cinq configurations, qui isole la
contribution de chaque brique (section~\ref{sec:ablation}).
\item La mise en évidence mesurée d'une dissociation entre performance sur la tâche
prétexte et qualité de la représentation apprise (section~\ref{sec:inversion}).
\item Un diagnostic sans réseau de neurones qui \emph{quantifie} le raccourci de couleur
responsable de cette dissociation : un simple histogramme de couleur résout la moitié de la
tâche prétexte lorsque les augmentations de couleur sont absentes
(section~\ref{sec:diagnostic}).
\item Une application web de démonstration projetant les représentations apprises en
2D, colorées par une classe réelle jamais vue à l'entraînement (section~\ref{sec:deploiement}).
\end{enumerate}

% =====================================================================
\section{Travaux liés}

\paragraph{SimCLR.} Chen \emph{et al.} \cite{chen2020simclr} formalisent le cadre que nous
reprenons : deux augmentations tirées de la même famille $\mathcal{T}$ produisent une paire
positive, un encodeur partagé $f(\cdot)$ en extrait des représentations, une tête de
projection $g(\cdot)$ les envoie dans l'espace où la perte contrastive s'applique, et
$g(\cdot)$ est jetée après le pré-entraînement. Leurs deux résultats les plus cités nous
concernent directement : la \emph{composition} d'augmentations compte davantage que chaque
augmentation prise isolément --- la paire recadrage $+$ distorsion de couleur étant
critique --- et la performance croît fortement avec la taille du lot, donc avec le nombre
de négatifs disponibles. Ce second point définit notre principale contrainte matérielle
(section~\ref{sec:limites}).

\paragraph{MoCo.} He \emph{et al.} \cite{he2020moco} attaquent exactement cette contrainte.
Plutôt que de tirer les négatifs du lot courant, MoCo maintient une \emph{file d'attente}
de représentations calculées aux itérations précédentes par un encodeur à moyenne mobile
(momentum). Le nombre de négatifs est ainsi découplé de la taille du lot que la mémoire GPU
autorise. C'est la référence naturelle pour tout projet contraint par un GPU unique.

\paragraph{Travaux complémentaires : se passer de négatifs.} Deux familles postérieures
suppriment le besoin de négatifs explicites, et donc la dépendance au lot. BYOL
\cite{grill2020byol} n'utilise que des paires positives : un réseau « en ligne » apprend à
prédire la sortie d'un réseau « cible » dont les poids sont une moyenne mobile des siens,
l'asymétrie entre les deux branches empêchant l'effondrement des représentations. DINO
\cite{caron2021dino} pousse cette idée sous forme de distillation auto-supervisée sans
étiquettes, en s'appuyant sur des transformeurs de vision, et fait émerger des cartes
d'attention qui segmentent les objets sans supervision.

\paragraph{Positionnement de notre travail.} Nous n'implémentons ni MoCo, ni BYOL, ni DINO :
le sujet demande un SimCLR simplifié, et notre budget de calcul est d'environ sept heures de
GPU. Ces travaux fixent en revanche l'horizon de nos résultats. Notre limite principale --- un
lot de 256, soit 510 négatifs par ancre, contre plusieurs milliers dans les
expériences de référence --- est précisément le problème que ces trois familles cherchent à
lever, chacune par un mécanisme différent. Nous y revenons en section~\ref{sec:limites}.

% =====================================================================
\section{Données}
\label{sec:donnees}

\subsection{Source et analyse exploratoire}

Nous utilisons CIFAR-10 \cite{krizhevsky2009cifar}, chargé via
\texttt{torchvision.datasets.CIFAR10} : 60\,000 images couleur de $32\times32$ pixels
réparties en 10 classes mutuellement exclusives (avion, automobile, oiseau, chat, cerf,
chien, grenouille, cheval, bateau, camion). La répartition officielle est de 50\,000 images
d'entraînement et 10\,000 images de test.

Trois propriétés du jeu de données ont guidé nos choix méthodologiques.

\begin{itemize}
\item \textbf{Équilibre parfait des classes.} Chaque classe compte exactement 6\,000
images, dont 5\,000 en entraînement et 1\,000 en test. L'accuracy est donc une métrique
légitime, sans qu'il soit nécessaire de recourir au F1 macro ou à une pondération ; le
hasard se situe à $10\,\%$, ce qui donne un point de repère commode.
\item \textbf{Statistiques des canaux.} Nous avons recalculé les moyennes et écarts-types
par canal directement sur les 50\,000 images d'entraînement, et ils reproduisent à la
quatrième décimale les constantes usuelles : moyennes $(0{,}4914;\ 0{,}4822;\ 0{,}4465)$,
écarts-types $(0{,}2470;\ 0{,}2435;\ 0{,}2616)$. Ces six nombres servent à la normalisation,
appliquée à l'identique dans toutes les expériences (\texttt{simclr/augmentations.py}). Ce
recalcul n'est pas une formalité : reprendre des constantes de normalisation trouvées en
ligne sans les vérifier est une source classique d'écart silencieux entre implémentations.
\item \textbf{Très petite résolution.} Une image CIFAR-10 fait $32\times32$ pixels, contre
$224\times224$ pour ImageNet. Cette différence d'un facteur 49 en nombre de pixels a deux
conséquences que nous assumons explicitement : elle impose d'adapter la tête d'un ResNet
conçu pour ImageNet (section~\ref{sec:encodeur}), et elle rend certaines augmentations
inopérantes, voire destructrices --- ce que notre ablation confirme pour le flou gaussien
(section~\ref{sec:ablation}).
\end{itemize}

\subsection{Ce que la couleur dit, et ce qu'elle ne dit pas}
\label{sec:eda-couleur}

Une quatrième propriété, moins souvent relevée, explique une partie des résultats de la
section~\ref{sec:ablation} et méritait d'être mesurée avant toute expérience.

Nous avons calculé, pour chaque classe, la couleur moyenne de ses images, puis deux
dispersions : celle des couleurs moyennes \emph{à l'intérieur} d'une classe, et celle des
couleurs moyennes \emph{entre} classes (figure~\ref{fig:edacouleur}). Le résultat est net.
L'écart-type intra-classe va de $27{,}6$ (camion) à $39{,}1$ (avion) niveaux, alors que
l'écart-type inter-classes ne vaut que $10{,}4$. La couleur varie donc environ trois fois
plus d'une image à l'autre \emph{dans} une classe qu'entre les classes elles-mêmes.

\fig{fig/fig_eda_couleur}{1.0}{Analyse exploratoire de la couleur sur les 50\,000 images
d'entraînement. À gauche, la couleur moyenne de chaque classe : neuf classes sur dix sont
des gris-bruns difficiles à distinguer, seuls l'avion et le bateau se détachent par leur
dominante bleue. À droite, la dispersion des couleurs moyennes à l'intérieur de chaque
classe, comparée à la dispersion entre classes (trait orange).}{fig:edacouleur}

\encadre{\textbf{Conséquence pour la suite.} La couleur d'une image CIFAR-10 est un
excellent identifiant \emph{d'instance} et un mauvais indice \emph{de catégorie}. Or la
tâche prétexte de SimCLR est précisément une tâche d'identification d'instance : elle
demande de reconnaître que deux vues viennent de la même image. La couleur est donc un
raccourci disponible, gratuit, et parfaitement suffisant pour résoudre la tâche prétexte
sans rien apprendre d'utile pour la classification. La section~\ref{sec:diagnostic} le
mesure directement.}

\subsection{Découpage, prévention des fuites}

Le jeu de test officiel ne sert qu'une fois, à la toute fin, pour les chiffres publiés :
aucun hyperparamètre n'a été choisi en le consultant. Pour disposer néanmoins d'un signal de
réglage, nous découpons le train officiel en une réserve étiquetable de 45\,000 images et
une validation de 5\,000 images (figure~\ref{fig:protocole}).

Ce découpage obéit à trois règles strictes.

\begin{enumerate}
\item Il est \textbf{stratifié} : chaque classe est représentée à parts égales dans la
validation.
\item Il utilise une \textbf{graine dédiée} ($7$), volontairement distincte des graines
d'expérience ($42$, $123$, $2026$), et n'est \textbf{jamais} refait. Si chaque graine
d'expérience voyait une validation différente, l'écart-type que nous rapportons mélangerait
deux sources de variation et ne serait plus interprétable.
\item Les trois fractions de labels sont \textbf{emboîtées} : les 450 images du régime
$1\,\%$ sont incluses dans les 4\,500 du régime $10\,\%$, elles-mêmes incluses dans les
45\,000. Ajouter des labels ne change donc jamais l'identité des labels déjà disponibles.
\end{enumerate}

Le script \texttt{scripts/inspecter\_split\_validation.py} vérifie automatiquement ces trois
propriétés (absence de recouvrement train/validation, stratification, emboîtement des
fractions) et refuse de valider le protocole si l'une d'elles est violée.

\fig{fig/fig_protocole}{1.0}{Protocole expérimental. Le pré-entraînement contrastif consomme
les 50\,000 images du train officiel \emph{sans leurs labels} ; les évaluations aval
n'utilisent que la réserve étiquetable, par fractions emboîtées ; le test officiel n'est
ouvert qu'une fois.}{fig:protocole}

\subsection{Choix de périmètre : CIFAR-10 seul}

Le sujet cite CIFAR-10 et STL-10 comme jeux de données possibles. Nous avons volontairement
restreint le périmètre à CIFAR-10. Aucune des exigences techniques du sujet ne mentionne
STL-10, et le budget de calcul disponible (une carte RTX 4060 Laptop, environ sept heures de
GPU au total) était mieux employé à consolider une seule comparaison --- trois graines par
point, quatre bras, cinq configurations d'ablation --- qu'à étaler deux comparaisons plus
fragiles sur deux jeux de données. C'est un choix de périmètre assumé, non un oubli ; nous
le rappelons en limites (section~\ref{sec:limites}).

% =====================================================================
\section{Méthode}
\label{sec:methode}

\subsection{Le cadre contrastif}

Le principe est illustré figure~\ref{fig:pipeline}. À partir d'une image $x$ tirée du jeu
non étiqueté, deux transformations $t, t' \sim \mathcal{T}$ sont tirées indépendamment et
produisent deux vues $\tilde{x}_i = t(x)$ et $\tilde{x}_j = t'(x)$ : c'est une \emph{paire
positive}. Un encodeur partagé $f(\cdot)$ en extrait les représentations
$h_i = f(\tilde{x}_i)$ et $h_j = f(\tilde{x}_j)$, puis une tête de projection $g(\cdot)$ les
envoie vers $z_i = g(h_i)$ et $z_j = g(h_j)$, où la perte contrastive s'applique. Toutes les
autres vues du lot constituent les \emph{paires négatives} : aucun échantillonnage explicite
de négatifs n'est nécessaire.

\fig{fig/fig_pipeline}{1.0}{Le cadre SimCLR tel que nous l'implémentons (schéma redessiné
d'après \cite{chen2020simclr}, figure 2). La perte s'applique sur $z$ ; c'est pourtant $h$
qui est réutilisé en aval, la tête de projection étant jetée après le pré-entraînement.}{fig:pipeline}

\subsection{Famille d'augmentations}

La famille $\mathcal{T}$ définit la tâche que le réseau doit résoudre : elle détermine à
quelles déformations la représentation devra être invariante. Notre configuration complète
enchaîne, dans cet ordre :

\begin{itemize}
\item \texttt{RandomResizedCrop(32, scale=(0.2, 1.0))} --- recadrage aléatoire couvrant de
$20\,\%$ à $100\,\%$ de la surface, puis redimensionnement à $32\times32$. Jamais retiré :
sans lui, les deux vues seraient géométriquement identiques et la tâche contrastive perdrait
son sens.
\item \texttt{RandomHorizontalFlip(p=0.5)} --- retournement horizontal.
\item \texttt{ColorJitter(0.4, 0.4, 0.4, 0.1)} appliqué avec $p = 0{,}8$ --- perturbation de
luminosité, contraste, saturation et teinte.
\item \texttt{RandomGrayscale(p=0.2)} --- conversion en niveaux de gris.
\item \texttt{GaussianBlur(kernel=3, sigma=(0.1, 1.0))} appliqué avec $p = 0{,}5$ ---
\emph{désactivé} dans la configuration principale, activé dans une variante d'ablation.
\item Conversion en tenseur et normalisation par les statistiques de la section~\ref{sec:donnees}.
\end{itemize}

Chaque brique est pilotable indépendamment, ce qui rend l'ablation de la
section~\ref{sec:ablation} possible sans dupliquer le code.

\subsection{Encodeur}
\label{sec:encodeur}

Nous utilisons un ResNet-18 \cite{he2016resnet}, adapté à la petite résolution de CIFAR-10
par deux modifications standard :

\begin{itemize}
\item la première convolution passe d'un noyau $7\times7$ de pas 2 à un noyau $3\times3$ de
pas 1. Sur une image de 32 pixels de côté, un noyau $7\times7$ de pas 2 réduit
immédiatement la carte à $16\times16$ et efface les détails fins ;
\item la couche de \emph{max pooling} qui suit est remplacée par l'identité, pour la même
raison : elle diviserait encore la résolution par deux avant tout traitement utile.
\end{itemize}

La couche entièrement connectée finale, qui produirait 1\,000 scores ImageNet, est également
remplacée par l'identité : la sortie de l'encodeur est la représentation
$h \in \mathbb{R}^{512}$ qui la précède. L'ensemble représente environ $11{,}2$ millions de
paramètres.

\subsection{Tête de projection}

$g(\cdot)$ est un perceptron à une couche cachée : $512 \rightarrow 512 \rightarrow 128$,
avec une activation ReLU intermédiaire. La perte s'applique donc sur $z \in \mathbb{R}^{128}$.
La tête est \textbf{jetée} après le pré-entraînement : toutes les évaluations aval
n'utilisent que $f(\cdot)$ et la représentation $h$.

\subsection{Perte NT-Xent}

Pour un lot de $N$ images, les $2N$ vues sont normalisées sur la sphère unité, puis toutes
les similarités cosinus sont calculées. Pour une ancre $i$ dont le positif est $j$ :

\begin{equation}
\mathcal{L}_i \;=\; -\log
\frac{\exp\!\big(\mathrm{sim}(z_i, z_j) / \tau\big)}
     {\sum_{k \neq i} \exp\!\big(\mathrm{sim}(z_i, z_k) / \tau\big)}
\qquad\text{avec}\qquad
\mathrm{sim}(u,v) = \frac{u^{\top} v}{\lVert u\rVert\,\lVert v\rVert}.
\label{eq:ntxent}
\end{equation}

Le dénominateur parcourt les $2N - 1$ autres vues du lot : à lot 256, cela fait 510 négatifs
par ancre. La température $\tau = 0{,}5$ contrôle la dureté de la discrimination : une
température basse force une séparation nette entre positifs et négatifs, une température
élevée lisse la distribution des similarités.

Notre implémentation (\texttt{simclr/contrastif.py}) exploite le fait que
l'équation~\eqref{eq:ntxent} est exactement une entropie croisée : après suppression de la
diagonale de la matrice de similarités, chaque ligne devient un vecteur de $2N-1$ scores
dont la cible est la position du positif, et \texttt{F.cross\_entropy} calcule directement
$\mathcal{L}$. Ce choix évite une boucle explicite et rend le code vérifiable ligne à ligne.

\subsection{Pré-entraînement}

Le pré-entraînement voit les 50\,000 images du train officiel, \emph{sans jamais lire un
seul label}. Hyperparamètres : 100 époques, lot 256, optimiseur AdamW avec taux
d'apprentissage $3\cdot10^{-4}$ et \emph{weight decay} $10^{-4}$, échauffement linéaire sur
10 époques suivi d'une décroissance en cosinus jusqu'à $10^{-6}$, température $0{,}5$,
précision mixte automatique, graine $42$. Un unique pré-entraînement est réalisé par
configuration d'augmentations, puis réutilisé pour tous les régimes de labels : c'est le
poste de calcul dominant et le recalculer à chaque point de la courbe serait un gaspillage.

\fig{fig/fig_preentrainement}{1.0}{Pré-entraînement contrastif, 100 époques. À gauche la
perte NT-Xent, à droite la part des ancres dont le positif est classé au premier rang parmi
les 511 candidats. Les deux courbes sont des diagnostics d'optimisation : elles indiquent
que l'entraînement se déroule normalement, mais \emph{ne mesurent pas} la qualité de la
représentation --- la section~\ref{sec:inversion} montre qu'elles peuvent même la
contredire.}{fig:preentrainement}

La durée médiane est de $30{,}4$ s par époque sur une RTX 4060 Laptop, soit environ 51
minutes pour 100 époques, avec un pic mémoire de $1{,}37$ Go. Quelques époques isolées
atteignent plusieurs minutes du fait de la contention de la machine ; nous en tirons la
conséquence en section~\ref{sec:limites} et ne rapportons \emph{aucune} comparaison de coût
en temps entre méthodes.

\subsection{Trois protocoles d'évaluation, plus un témoin}
\label{sec:protocoles}

Un encodeur pré-entraîné peut être exploité de deux manières, et le sujet demande de
comparer avec et sans pré-entraînement. Nous évaluons donc quatre bras.

\begin{enumerate}
\item \textbf{Évaluation linéaire} (\emph{linear probe}). L'encodeur est gelé ; seule une
couche $\mathrm{Linear}(512, 10)$ est entraînée sur les représentations $h$ standardisées.
Environ 5\,130 paramètres entraînables. C'est le protocole standard pour comparer la qualité
\emph{intrinsèque} de deux pré-entraînements, et c'est celui que nous utilisons dans
l'ablation.
\item \textbf{Fine-tuning complet.} L'encodeur pré-entraîné et la tête de classification
s'adaptent ensemble. Environ $11{,}2$ millions de paramètres entraînables. C'est la
performance réellement utile en pratique.
\item \textbf{Supervisé \emph{from scratch}.} Strictement le même entraînement que le
fine-tuning, à partir de poids aléatoires. C'est la baseline concurrente.
\item \textbf{Témoin à encodeur aléatoire gelé.} Le protocole d'évaluation linéaire appliqué
à un ResNet-18 jamais entraîné. Ce bras n'est demandé nulle part ; nous l'ajoutons parce
qu'il est le seul moyen de distinguer ce que le pré-entraînement apporte de ce que
l'architecture et le protocole d'évaluation apportent à eux seuls.
\end{enumerate}

\subsection{Égalité des budgets de calcul}

Le piège le plus courant dans ce type de comparaison consiste à opposer deux méthodes qui
n'ont pas reçu le même budget d'entraînement. Le tableau~\ref{tab:hyper} liste les
hyperparamètres des deux bras comparés, extraits des fichiers
\texttt{resultats.json} produits par les runs eux-mêmes, dans
\texttt{outputs/finetuning\_simclr\_45k/} et
\texttt{outputs/supervise\_from\_scratch\_45k/}.

\begin{table}[htbp]
\centering\small
\caption{Hyperparamètres des deux bras comparés. La seule différence entre les deux colonnes
est l'initialisation des poids de l'encodeur.}
\label{tab:hyper}
\begin{tabularx}{\textwidth}{@{}lXX@{}}
\toprule
& \textbf{Fine-tuning SimCLR} & \textbf{Supervisé \emph{from scratch}} \\
\midrule
Initialisation de l'encodeur & poids pré-entraînés (100 époques SSL) & poids aléatoires \\
Architecture & ResNet-18 adapté, $\sim$11,2 M & ResNet-18 adapté, $\sim$11,2 M \\
Époques & 100 & 100 \\
Taille de lot & 256 & 256 \\
Optimiseur & SGD, momentum 0,9, Nesterov & SGD, momentum 0,9, Nesterov \\
Taux d'apprentissage & 0,03 (cosinus $\rightarrow 10^{-4}$) & 0,03 (cosinus $\rightarrow 10^{-4}$) \\
\emph{Weight decay} & $5\cdot10^{-4}$ & $5\cdot10^{-4}$ \\
Augmentations d'entraînement & \texttt{RandomCrop(32, pad=4)} $+$ \texttt{Flip} & \texttt{RandomCrop(32, pad=4)} $+$ \texttt{Flip} \\
Précision mixte & activée & activée \\
Découpage / graines & \texttt{seed\_split} 7 ; 42, 123, 2026 & \texttt{seed\_split} 7 ; 42, 123, 2026 \\
\bottomrule
\end{tabularx}
\end{table}

\subsection{Choix du taux d'apprentissage}

Le taux d'apprentissage est le seul hyperparamètre que nous avons réglé, et nous l'avons
réglé \emph{symétriquement} : la même grille $\{0{,}2;\ 0{,}1;\ 0{,}03;\ 0{,}01;\ 0{,}003\}$
a été balayée pour les deux méthodes, sur la validation, au régime $10\,\%$ des labels
(figure~\ref{fig:lr}). Aucune des deux ne bénéficie donc d'un avantage de réglage.

Les deux courbes atteignent leur maximum en $0{,}03$, valeur retenue pour les deux bras.
La valeur $0{,}2$ dégrade fortement les deux méthodes, ce qui indique que la grille encadre
bien l'optimum plutôt que de s'arrêter à sa frontière.

\fig{fig/fig_lr}{0.78}{Balayage symétrique du taux d'apprentissage sur la validation, au
régime 10\,\% des labels. Même grille pour les deux méthodes.}{fig:lr}

\subsection{Métriques, graines et reproductibilité}

La métrique publiée est l'accuracy top-1 sur les 10\,000 images du test officiel. Chaque
cellule des tableaux est la moyenne de trois répétitions (graines $42$, $123$, $2026$),
accompagnée de l'écart-type. Les graines contrôlent l'initialisation, l'ordre des lots et le
tirage des sous-ensembles étiquetés ; le découpage train/validation, lui, reste fixe.
Chaque ligne de résultat est reproductible par une commande unique, documentée dans
\texttt{docs/REPRODUCTION.md} du dépôt.

% =====================================================================
\section{Étude d'ablation : quelles augmentations comptent ?}
\label{sec:ablation}

\subsection{Question et protocole}

Le sujet demande d'étudier l'effet du choix des augmentations sur la qualité des
représentations apprises. Nous adoptons un design \textbf{additif} : on part du recadrage
seul et on ajoute \emph{une} brique à la fois. C'est la règle d'or d'une ablation --- une
seule variable change entre deux lignes, tout le reste est identique.

Les cinq configurations partagent rigoureusement : ResNet-18 adapté, 100 époques, lot 256,
AdamW $3\cdot10^{-4}$, échauffement 10 époques puis cosinus, température $0{,}5$, précision
mixte, graine de pré-entraînement $42$. L'évaluation est un \emph{linear probe} sur encodeur
gelé, moyenné sur trois graines.

Nous rapportons cette évaluation à \textbf{deux} régimes d'annotation. À $10\,\%$ des labels
d'abord, parce que c'est le protocole le plus discriminant : avec peu d'étiquettes, le
classifieur ne peut pas compenser une représentation médiocre, et la qualité de l'encodeur
pèse donc au maximum. À $100\,\%$ ensuite, parce que la mesure y est nettement moins
bruitée --- l'écart-type de la configuration la plus instable passe de $\pm 0{,}54$ à
$\pm 0{,}06$. Les deux commandes sont \texttt{python -m scripts.reproduire\_ablation} et la
même avec l'option \texttt{-{}-fraction-probe 0.10} ; ensemble, elles régénèrent
l'intégralité du tableau~\ref{tab:ablation}.

\subsection{Résultats}

\begin{table}[htbp]
\centering\small
\caption{Ablation additive des augmentations. Les deux colonnes de gauche sont des
\emph{diagnostics} de la tâche prétexte et ne servent jamais au classement. Les deux
colonnes de droite donnent l'accuracy de l'évaluation linéaire aux deux régimes
d'annotation ; le classement est identique dans les deux. Écart-type sur trois graines
de \emph{probe}.}
\label{tab:ablation}
\begin{tabular}{@{}lcccc@{}}
\toprule
& \multicolumn{2}{c}{\emph{diagnostics de la tâche prétexte}} & \multicolumn{2}{c}{\textbf{accuracy de l'évaluation linéaire}} \\
\cmidrule(lr){2-3}\cmidrule(lr){4-5}
\textbf{Configuration} & \textbf{Perte NT-Xent} & \textbf{Positif 1\textsuperscript{er}} & \textbf{à 10\,\% des labels} & \textbf{à 100\,\% des labels} \\
\midrule
crop seul                    & 4,3739 & 98,01\,\% & $53{,}29 \pm 0{,}54$ & $59{,}99 \pm 0{,}06$ \\
\quad $+$ flip               & 4,3802 & 97,95\,\% & $52{,}68 \pm 0{,}40$ & $59{,}33 \pm 0{,}11$ \\
\quad $+$ color jitter       & 4,4674 & 91,68\,\% & $68{,}16 \pm 0{,}39$ & $72{,}79 \pm 0{,}04$ \\
\quad $+$ niveaux de gris \textsuperscript{(a)} & 4,5502 & 80,47\,\% & $\mathbf{78{,}83 \pm 0{,}22}$ & $\mathbf{81{,}87 \pm 0{,}04}$ \\
\quad $+$ flou gaussien      & 4,5601 & 78,86\,\% & $78{,}84 \pm 0{,}13$ & $81{,}56 \pm 0{,}03$ \\
\bottomrule
\end{tabular}

\vspace{3pt}
{\footnotesize\textsuperscript{(a)} Configuration retenue pour toutes les expériences de la
section~\ref{sec:resultats}.}
\end{table}

\fig{fig/fig_ablation}{1.0}{Ablation des augmentations. À gauche, le niveau atteint par
chaque configuration ; à droite, la contribution propre de la brique ajoutée.}{fig:ablation}

\subsection{Discussion, variante par variante}

La figure~\ref{fig:ablation} décompose ces chiffres : à gauche le niveau atteint par chaque
configuration, à droite la contribution propre de la brique ajoutée. Nous commentons chaque
variante ci-dessous.

\paragraph{Les deux régimes donnent le même classement.} C'est le premier point à relever :
l'ordre des cinq configurations est rigoureusement identique à $10\,\%$ et à $100\,\%$ des
labels. La conclusion de l'ablation ne dépend donc pas du régime d'annotation choisi pour
l'évaluer, ce qui la rend nettement plus robuste qu'une mesure à un seul point.

Les \emph{amplitudes}, en revanche, diffèrent, et dans un sens instructif. À $10\,\%$ des
labels, les deux briques de couleur apportent $+26{,}16$ points ($+15{,}48$ pour le color
jitter, $+10{,}68$ pour les niveaux de gris) contre $+22{,}54$ à $100\,\%$. Autrement dit :
plus les étiquettes sont rares, plus la qualité de la représentation compte. C'est cohérent
avec le résultat central du projet --- quand les labels abondent, le classifieur peut
compenser une représentation médiocre ; quand ils manquent, il ne le peut plus.

\paragraph{Le retournement horizontal : $-0{,}66$ point à $100\,\%$, $-0{,}62$ à $10\,\%$.}
Aucun effet mesurable. Le
retournement ne modifie ni les couleurs ni le contenu de l'image ; il n'ajoute donc aucune
difficulté à la tâche contrastive. L'écart observé est du même ordre que le bruit
d'évaluation et nous ne le rapportons \emph{pas} comme un effet négatif.

\paragraph{La distorsion de couleur : $+13{,}46$ points à $100\,\%$, $+15{,}48$ à $10\,\%$.}
C'est le gain le plus important de l'étude, et de loin. Il retrouve indépendamment le résultat central de
\cite{chen2020simclr} : sans perturbation des couleurs, les deux vues d'une même image
conservent des histogrammes de couleur quasi identiques, ce qui offre au réseau un raccourci
lui permettant de résoudre la tâche sans jamais analyser la forme.

\paragraph{Les niveaux de gris : $+9{,}08$ points à $100\,\%$, $+10{,}68$ à $10\,\%$.}
Second gain majeur. Cette brique
supprime le résidu du raccourci précédent : lorsqu'une des deux vues est convertie en
niveaux de gris, la comparaison par teinte devient impossible et la forme redevient le seul
signal exploitable. Cumulées, les deux briques de couleur apportent $+22{,}54$ points ---
soit la quasi-totalité de l'écart entre la pire et la meilleure configuration.

\paragraph{Le flou gaussien : $-0{,}31$ point à $100\,\%$, $+0{,}01$ à $10\,\%$.} Aucun
effet mesurable non plus --- et le fait que le signe s'inverse entre les deux régimes, pour
une amplitude aussi faible, confirme qu'il s'agit de bruit et non d'un effet. Ce résultat
diffère de celui rapporté sur ImageNet, où le flou aide. La divergence s'explique par la
résolution et non par une contradiction : sur une image de $224\times224$ pixels, un noyau
de flou reste local ; sur une image de $32\times32$, un noyau $3\times3$ couvre presque
$10\,\%$ de la largeur de l'image et détruit une part significative de l'information
disponible. Nous en concluons que le flou n'est pas transférable tel quel à cette échelle,
non qu'il serait inutile en général.

\encadre{\textbf{Réponse à la question du sujet.} Sur CIFAR-10, ce sont les augmentations de
\emph{couleur} qui portent la qualité de la représentation : $+22{,}54$ points à $100\,\%$
des labels et $+26{,}16$ points à $10\,\%$ pour la distorsion de couleur et les niveaux de
gris réunis, contre des écarts non mesurables pour le retournement horizontal et le flou
gaussien. Le classement est identique aux deux régimes.}

\subsection{Bien résoudre la tâche prétexte n'est pas bien apprendre}
\label{sec:inversion}

Le résultat le plus instructif de cette ablation n'est pas le classement lui-même, mais sa
comparaison avec le classement par la perte.

La configuration \texttt{crop seul} obtient \emph{la meilleure} perte NT-Xent finale
($4{,}3739$) et place $98{,}01\,\%$ des positifs au premier rang. Elle résout donc presque
parfaitement la tâche prétexte --- et produit la \emph{pire} représentation de l'étude
($59{,}99\,\%$). Symétriquement, la configuration complète obtient une perte plus élevée
($4{,}5502$), un taux de positifs au premier rang bien plus faible ($80{,}47\,\%$), et la
meilleure représentation ($81{,}87\,\%$). La figure~\ref{fig:inversion} montre que le
classement s'inverse presque exactement d'un critère à l'autre.

\fig{fig/fig_inversion}{0.92}{Rang de chaque configuration selon la perte NT-Xent (à gauche)
et selon l'accuracy de l'évaluation linéaire (à droite).}{fig:inversion}

Le mécanisme est identifiable, et c'est ce qui rend l'observation solide
(figure~\ref{fig:raccourci}) : retirer les perturbations de couleur laisse aux deux vues le
même histogramme de couleur, ce qui permet de les apparier sans analyser la forme. La tâche
devient facile, la perte descend --- et la représentation apprise se réduit à un descripteur
de couleur, inutile en aval.

\fig{fig/fig_raccourci}{1.0}{Le raccourci de couleur, mécanisme derrière l'inversion.}{fig:raccourci}

\encadre{\textbf{À retenir.} Une perte contrastive plus basse peut signaler une tâche
prétexte trop facile, donc une représentation plus pauvre. Le seul critère de classement
valide est la performance sur la tâche aval. C'est pourquoi la commande de reproduction de
l'ablation affiche explicitement les deux colonnes de diagnostic comme
\emph{non décisionnelles}.}

Nous nous gardons de transformer cette observation en loi générale. La corrélation de rang
entre perte et accuracy vaut $\rho = +0{,}80$, mais sur $n = 5$ configurations seulement,
ce qui ne permet aucune conclusion statistique ($p \approx 0{,}10$). Ce que nous affirmons
est plus modeste et mieux étayé : sur nos cinq configurations, les deux classements
s'opposent, et un mécanisme concret explique pourquoi.

\subsection{Le raccourci, mesuré sans aucun réseau de neurones}
\label{sec:diagnostic}

L'explication précédente reste une interprétation tant qu'elle n'est pas mesurée. Nous avons
donc conçu un diagnostic qui la teste directement, et dont la particularité est de
\textbf{n'employer aucun réseau de neurones} : il ne mesure pas ce que notre modèle a appris,
mais ce que les données et les augmentations rendent possible.

\paragraph{Protocole.} Sur 2\,000 images d'entraînement, nous produisons deux vues par image
avec la chaîne d'augmentations étudiée. Chaque vue est résumée par un histogramme de couleur
de $3 \times 8 = 24$ bins, normalisé --- une description qui ignore complètement la forme, la
texture et la position. Nous mesurons ensuite deux quantités : le taux de vues dont le
véritable positif est le plus proche voisin au sens de la distance $L^1$ entre histogrammes,
parmi les 2\,000 candidates ; et le rapport entre la distance moyenne aux négatifs et la
distance moyenne au positif. Le hasard se situe à $1/2000 = 0{,}05\,\%$.

\begin{table}[htbp]
\centering\small
\caption{Diagnostic du raccourci de couleur, mené sur les données seules. Un simple
histogramme de couleur suffit-il à résoudre la tâche prétexte ?}
\label{tab:diagnostic}
\begin{tabular}{@{}lccc@{}}
\toprule
\textbf{Augmentations appliquées} & \textbf{Positif retrouvé} & \textbf{Distance nég. / pos.}
& \textbf{Accuracy de l'encodeur} \\
\midrule
crop seul                              & $\mathbf{51{,}9\,\%}$ & $\times 3{,}02$ & $59{,}99\,\%$ \\
\quad $+$ color jitter                 & $6{,}1\,\%$           & $\times 1{,}35$ & $72{,}79\,\%$ \\
\quad $+$ niveaux de gris              & $3{,}8\,\%$           & $\times 1{,}31$ & $\mathbf{81{,}87\,\%}$ \\
\midrule
\emph{hasard}                          & $0{,}05\,\%$          & $\times 1{,}00$ & $10\,\%$ \\
\bottomrule
\end{tabular}
\end{table}

\fig{fig/fig_diagnostic_couleur}{1.0}{Le raccourci de couleur, mesuré directement sur les
données. À gauche, la part des vues dont le vrai positif est retrouvé au premier rang par
la seule couleur, parmi 2\,000 candidates. À droite, le rapport entre distance aux négatifs
et distance au positif : au-dessus de $\times 1$, la couleur porte un signal
exploitable.}{fig:diagnostic}

\paragraph{Résultat.} Avec le recadrage seul, un histogramme de couleur retrouve la bonne
image dans $51{,}9\,\%$ des cas parmi 2\,000 candidates, soit plus de mille fois le hasard.
La tâche prétexte est donc, littéralement, à moitié résolue avant même qu'un réseau ne
regarde une seule forme. Ajouter la distorsion de couleur fait tomber ce taux à $6{,}1\,\%$,
et les niveaux de gris à $3{,}8\,\%$ ; le rapport de distances passe de $\times 3{,}02$ à
$\times 1{,}31$, c'est-à-dire presque plus rien.

\paragraph{Interprétation.} Ce diagnostic transforme notre explication en mesure, et il
éclaire les trois colonnes du tableau~\ref{tab:ablation} d'un seul coup. Les $98{,}01\,\%$ de
positifs au premier rang obtenus par la configuration \texttt{crop seul} ne témoignent pas
d'une représentation supérieure : plus de la moitié de ce score était accessible à un
histogramme de couleur, sans le moindre apprentissage. Inversement, la chute du taux de
positifs à $80{,}47\,\%$ pour la configuration complète n'est pas une dégradation --- c'est
le signe que le raccourci a été fermé et que le réseau a dû apprendre autre chose. Le gain
de $+22{,}54$ points en aval est le prix de cette fermeture.

\encadre{\textbf{Ce que nous pouvons affirmer.} Le raccourci de couleur n'est pas une
hypothèse commode pour expliquer nos chiffres : il est mesurable sur les données brutes,
il est massif ($51{,}9\,\%$ contre $0{,}05\,\%$ au hasard), et les deux augmentations qui le
neutralisent sont exactement celles qui apportent les $+22{,}54$ points de notre ablation.}

\paragraph{Limites de ce diagnostic.} Il mesure une \emph{borne inférieure} de ce que la
couleur permet, pas une borne supérieure de ce que le réseau exploite : un descripteur de
couleur plus fin (histogrammes spatiaux, moments d'ordre supérieur) ferait probablement
mieux. Il porte sur 2\,000 images d'un seul lot de données, avec une graine unique. Enfin,
notre implémentation des transformations de couleur est une reproduction en \texttt{numpy}
de \texttt{ColorJitter} et \texttt{RandomGrayscale} ; elle en respecte les paramètres et
l'ordre d'application, mais n'est pas le code de \texttt{torchvision} lui-même.

\subsection{Ce que cette ablation ne permet pas d'affirmer}

Trois limites doivent accompagner le tableau~\ref{tab:ablation}, et nous préférons les
énoncer plutôt que d'attendre qu'elles nous soient opposées.

\begin{itemize}
\item Les $\pm 0{,}0X$ sont l'écart-type sur les \textbf{trois graines du linear probe
uniquement}. Le pré-entraînement n'a tourné qu'\textbf{une seule fois} par configuration
(graine $42$) : la variabilité d'un pré-entraînement à l'autre est inconnue. Trois graines de
pré-entraînement sur cinq configurations représenteraient environ dix heures de GPU
supplémentaires, hors de notre budget.
\item Par conséquent, $+13{,}46$ et $+9{,}08$ sont des effets trop grands pour être du bruit
et nous les traitons comme des conclusions solides ; $-0{,}66$ et $-0{,}31$ sont rapportés
comme « aucun effet mesurable » et jamais comme des effets négatifs.
\item L'ablation est additive, non factorielle : elle mesure l'apport de chaque brique
\emph{dans l'ordre où nous les ajoutons}. Elle ne dit rien des interactions entre briques
prises dans un autre ordre --- par exemple ce que donnerait « crop $+$ niveaux de gris » sans
distorsion de couleur.
\end{itemize}

% =====================================================================
\section{Résultats}
\label{sec:resultats}

\subsection{Tableau principal}

Le tableau~\ref{tab:principal} et la figure~\ref{fig:courbe} constituent le résultat central
du projet.

\begin{table}[htbp]
\centering\small
\caption{Accuracy top-1 sur le test officiel (10\,000 images), en \%, moyenne $\pm$
écart-type sur trois graines. Toutes les méthodes sont appariées : même découpage, mêmes
graines, mêmes sous-ensembles étiquetés.}
\label{tab:principal}
\begin{tabular}{@{}llccc@{}}
\toprule
\textbf{Méthode} & \textbf{Poids entraînés} & \textbf{450 (1\,\%)} & \textbf{4\,500 (10\,\%)} & \textbf{45\,000 (100\,\%)} \\
\midrule
Encodeur aléatoire gelé \textsuperscript{(t)} & $\sim$5\,130   & $25{,}66 \pm 1{,}18$ & $33{,}25 \pm 0{,}46$ & $42{,}30 \pm 0{,}02$ \\
Linear probe SimCLR      & $\sim$5\,130   & $72{,}00 \pm 0{,}40$ & $78{,}83 \pm 0{,}22$ & $81{,}87 \pm 0{,}04$ \\
Fine-tuning SimCLR       & $\sim$11,2 M   & $\mathbf{72{,}10 \pm 0{,}86}$ & $\mathbf{82{,}71 \pm 0{,}35}$ & $\mathbf{93{,}56 \pm 0{,}09}$ \\
Supervisé \emph{from scratch} & $\sim$11,2 M & $39{,}52 \pm 0{,}58$ & $70{,}25 \pm 0{,}56$ & $93{,}13 \pm 0{,}08$ \\
\midrule
\multicolumn{2}{@{}l}{Écart fine-tuning $-$ supervisé} & $+32{,}58$ & $+12{,}46$ & $+0{,}43$ \\
\bottomrule
\end{tabular}

\vspace{3pt}
{\footnotesize\textsuperscript{(t)} Témoin diagnostic : même architecture et même protocole
d'évaluation, poids jamais entraînés.}
\end{table}

\fig{fig/fig_accuracy_labels}{1.0}{Accuracy de test en fonction du nombre d'images
étiquetées, avec et sans pré-entraînement auto-supervisé. Les barres d'erreur représentent
l'écart-type sur trois graines ; elles sont souvent plus petites que les marqueurs.}{fig:courbe}

\subsection{Lecture des trois régimes}

\paragraph{450 labels : le pré-entraînement change la nature du problème.} L'écart de
$+32{,}58$ points est le résultat le plus net de l'étude. Avec 45 images étiquetées par
classe, la baseline supervisée atteint $39{,}52\,\%$ ; le même réseau, initialisé par le
pré-entraînement contrastif, atteint $72{,}10\,\%$. Formulé du point de vue du coût
d'annotation : le modèle pré-entraîné atteint $72{,}1\,\%$ avec 450 labels, tandis que la
baseline a besoin de 4\,500 labels --- dix fois plus --- pour n'atteindre que $70{,}25\,\%$.
L'économie d'annotation dépasse $90\,\%$ à performance équivalente.

\paragraph{450 labels, second constat : dégeler ne sert à rien.} À ce régime, le fine-tuning
complet ($72{,}10\,\%$) et l'évaluation linéaire sur encodeur gelé ($72{,}00\,\%$) sont
séparés de $0{,}10$ point, soit nettement moins que l'écart-type du premier ($\pm 0{,}86$).
Autrement dit : avec très peu de labels, entraîner $11{,}2$ millions de paramètres n'apporte
rien de mesurable par rapport à en entraîner $5\,130$. Le gradient disponible est trop
maigre pour améliorer l'encodeur, et le geler évite qu'il ne se dégrade. C'est un résultat
directement actionnable en pratique : à faible supervision, l'évaluation linéaire est aussi
bonne et deux mille fois moins coûteuse en paramètres.

\paragraph{4\,500 labels : le régime intermédiaire.} L'écart se réduit à $+12{,}46$ points.
Le fine-tuning dépasse ici l'évaluation linéaire ($82{,}71$ contre $78{,}83$) : il y a
désormais assez de signal pour que l'adaptation de l'encodeur devienne rentable.

\paragraph{45\,000 labels : parité.} Avec toutes les annotations, les deux méthodes se
rejoignent ($93{,}56$ contre $93{,}13$). Le bénéfice du SSL s'annule quand les labels
abondent, exactement comme l'anticipe la littérature. Nous insistons sur un point : nous
parlons de \textbf{parité} et nous ne revendiquons \emph{pas} l'écart de $+0{,}43$ point
comme une supériorité --- voir section~\ref{sec:limites}.

\paragraph{Le témoin.} Sans lui, un lecteur sceptique pourrait attribuer les 32 points au
protocole d'évaluation plutôt qu'au pré-entraînement. Le témoin ferme cette porte : à 450
labels, un encodeur aléatoire gelé évalué exactement comme le nôtre atteint $25{,}66\,\%$,
contre $72{,}00\,\%$ pour l'encodeur pré-entraîné --- $46{,}34$ points d'écart, à
architecture, protocole et nombre de paramètres entraînables strictement identiques. La seule
variable qui change est le contenu des poids de l'encodeur.

Le témoin apporte un second enseignement, moins attendu. Un encodeur \emph{jamais entraîné}
atteint tout de même $42{,}30\,\%$ avec 45\,000 labels, très au-dessus des $10\,\%$ du
hasard. Une convolution aléatoire n'est donc pas un bruit : la structure convolutionnelle
elle-même --- localité, partage de poids, invariance par translation --- extrait déjà des
descripteurs partiellement linéairement séparables. C'est un rappel utile pour lire les
comparaisons de la littérature : la bonne référence n'est jamais le hasard à $10\,\%$, mais
ce que la même architecture obtient sans apprentissage.

\subsection{Lecture en coût d'annotation}
\label{sec:cout-annotation}

Les points d'accuracy sont la métrique naturelle du chercheur ; le nombre d'annotations est
la métrique naturelle de celui qui finance le projet. La courbe se lit aussi horizontalement.

\begin{table}[htbp]
\centering\small
\caption{Lecture horizontale de la figure~\ref{fig:courbe} : combien d'images faut-il
annoter pour atteindre un niveau donné ?}
\label{tab:cout}
\begin{tabular}{@{}lccc@{}}
\toprule
\textbf{Niveau visé} & \textbf{Avec pré-entraînement} & \textbf{Sans pré-entraînement}
& \textbf{Économie} \\
\midrule
$\approx 70\,\%$ & 450 images ($72{,}10\,\%$)   & 4\,500 images ($70{,}25\,\%$)  & $\times 10$ \\
$\approx 82\,\%$ & 4\,500 images ($82{,}71\,\%$) & entre 4\,500 et 45\,000       & $> \times 3$ \\
$\approx 93\,\%$ & 45\,000 images               & 45\,000 images                 & aucune \\
\bottomrule
\end{tabular}
\end{table}

La première ligne est celle qui porte l'argument économique : atteindre $70\,\%$ demande
450 annotations avec pré-entraînement contre 4\,500 sans, soit une réduction de $90\,\%$ du
volume à annoter. Si l'on transpose à un contexte où chaque annotation demande l'avis d'un
spécialiste, c'est la différence entre une demi-journée de travail et une semaine.

La deuxième ligne est encadrée plutôt que chiffrée, et il faut le dire : nous n'avons mesuré
que trois points, nous ne pouvons donc pas interpoler le nombre exact d'annotations que la
baseline supervisée exigerait pour atteindre $82\,\%$. Nous savons seulement qu'il est
supérieur à 4\,500 (elle y plafonne à $70{,}25\,\%$) et inférieur à 45\,000 (elle y atteint
$93{,}13\,\%$). Une mesure à 20\,000 labels lèverait cette incertitude ; elle n'était pas
dans notre budget.

\subsection{Le régime de faible supervision, dans le détail}

Un point mérite d'être isolé car il est directement actionnable et n'apparaît dans aucun
des documents de référence du sujet.

À 450 labels, le fine-tuning complet ($72{,}10 \pm 0{,}86$) et l'évaluation linéaire sur
encodeur gelé ($72{,}00 \pm 0{,}40$) sont séparés de $0{,}10$ point, soit huit fois moins
que l'écart-type du premier. Entraîner $11{,}2$ millions de paramètres n'apporte donc
\emph{rien de mesurable} par rapport à en entraîner $5\,130$ --- un rapport de 1 à 2\,180.
Deux mécanismes se compensent : avec 45 exemples par classe, le gradient est trop pauvre
pour améliorer un encodeur déjà bon, et dégeler l'encodeur l'expose à une dégradation par
sur-ajustement, que l'accuracy d'entraînement à $100\,\%$ (section~\ref{sec:surapp})
confirme.

Cet écart se creuse ensuite : $+3{,}88$ points en faveur du fine-tuning à 4\,500 labels,
$+11{,}69$ à 45\,000. Le seuil où le dégel devient rentable se situe donc entre 450 et
4\,500 annotations dans notre configuration. En pratique, cela signifie qu'un praticien
disposant de très peu d'étiquettes a intérêt à geler son encodeur : c'est aussi précis, bien
plus rapide, et le modèle tient dans quelques kilo-octets une fois l'encodeur partagé.

\subsection{Comparaison à la littérature}

Les implémentations de référence de SimCLR rapportent, sur CIFAR-10, des évaluations
linéaires de l'ordre de $94\,\%$, avec un encodeur plus large, des lots de plusieurs milliers
d'images et des pré-entraînements de l'ordre de 1\,000 époques \cite{chen2020simclr}. Notre
évaluation linéaire plafonne à $81{,}87\,\%$. L'écart n'a rien de surprenant et se lit comme
un écart de moyens : nous utilisons un ResNet-18 et non un ResNet plus profond ou plus large,
un lot de 256 et non de plusieurs milliers, et 100 époques et non 1\,000, soit un dixième du
budget de pré-entraînement. Ce qui se transfère malgré cette échelle réduite, ce n'est pas le
chiffre absolu mais la \emph{structure} du résultat : l'écart maximal en régime de faible
supervision, sa disparition à supervision complète, et la primauté des augmentations de
couleur. Ces trois conclusions qualitatives sont reproduites.

\subsection{Analyse critique : à quel coût ?}

La question « le modèle est-il réellement meilleur, et à quel prix ? » appelle une réponse
nuancée.

\begin{itemize}
\item \textbf{En régime de faible supervision, oui, très largement}, et le coût
supplémentaire est de 100 époques de pré-entraînement --- payées \emph{une seule fois} et
amorties sur les trois régimes de labels, puisque le même encodeur est réutilisé partout.
\item \textbf{En régime de supervision complète, non.} Les deux méthodes sont à parité et le
pré-entraînement représente alors un coût de calcul pur, sans contrepartie mesurable.
\item \textbf{Nous ne rapportons aucune comparaison de temps d'exécution.} Les mesures de
durée de nos runs sont inexploitables : la machine était partagée avec d'autres charges, et
une même configuration à $100\,\%$ des labels a mis $2\,104$ s pour la graine $123$ et
$11\,914$ s pour la graine $2026$. Publier un ratio de coût à partir de ces chiffres serait
malhonnête ; nous préférons ne rien affirmer que d'affirmer une mesure que le protocole ne
soutient pas.
\end{itemize}

\subsection{Surapprentissage}
\label{sec:surapp}

Dans les neuf configurations supervisées, l'accuracy d'entraînement atteint $100\,\%$, y
compris au régime de 450 labels où l'accuracy de test est de $39{,}52\,\%$. Le réseau mémorise
donc intégralement les 450 exemples disponibles sans en généraliser grand-chose --- un écart
de plus de 60 points entre train et test. C'est précisément l'espace que le pré-entraînement
vient combler : il fournit à l'encodeur une structure apprise sur 50\,000 images, que les
450 labels n'ont plus qu'à réorienter vers la tâche de classification.

Ce constat a une conséquence méthodologique qu'il vaut mieux énoncer : à 450 labels, la
baseline supervisée n'est pas limitée par sa capacité mais par son information. Augmenter la
taille du modèle ne l'aiderait pas ; ajouter de la régularisation déplacerait le problème
sans le résoudre. Le pré-entraînement auto-supervisé agit sur la seule variable qui manque,
l'information --- et il la prend là où elle est gratuite, dans les 50\,000 images non
étiquetées.

% =====================================================================
\section{Déploiement : application de démonstration}
\label{sec:deploiement}

Conformément à l'exigence allégée du sujet, l'application sert à \emph{illustrer} les
résultats, pas à constituer un produit fini. Elle prend la forme d'une page web autonome
(\texttt{demo/index.html}) : un unique fichier HTML, sans serveur, sans dépendance réseau et
sans conteneur Docker, qui s'ouvre directement dans un navigateur. Le script
\texttt{scripts/preparer\_demo.py} extrait les représentations, calcule les projections et
embarque le tout dans la page.

La démonstration répond à la demande du sujet --- « visualisation des représentations
apprises, projection t-SNE/UMAP colorée par classe réelle, non utilisée à l'entraînement » ---
et permet de :

\begin{itemize}
\item comparer côte à côte les projections t-SNE et PCA de \textbf{six} encodeurs :
l'encodeur aléatoire et les cinq configurations d'ablation ;
\item colorer les 2\,000 points par la classe réelle, qui n'a jamais été vue au
pré-entraînement, et allumer ou éteindre chaque classe ;
\item sélectionner une image et afficher ses plus proches voisins dans l'espace des
représentations, ce qui rend le raccourci de couleur visible à l'\oe il nu sur la
configuration \texttt{crop seul} ;
\item consulter la matrice de confusion de l'évaluation linéaire.
\end{itemize}

\fig{fig/fig_tsne}{1.0}{Projection t-SNE des représentations $h \in \mathbb{R}^{512}$ de
2\,000 images de test, telle que la démonstration l'affiche. Les couleurs correspondent aux
classes réelles, jamais utilisées au pré-entraînement ; les étiquettes sont placées au
centroïde de chaque amas.}{fig:tsne}

La figure~\ref{fig:tsne} résume ce que la démonstration donne à voir. L'encodeur aléatoire
produit un nuage informe où les couleurs se mélangent. Le pré-entraînement au recadrage seul
crée une structure, mais sans séparation des classes. Les augmentations complètes font
apparaître des amas nets, alignés sur les classes réelles alors qu'aucun label n'a été
utilisé : c'est la preuve visuelle que la tâche contrastive a capté une structure
sémantique, et non un artefact de bas niveau.

% =====================================================================
\section{Discussion et limites}
\label{sec:limites}

\subsection{Limite principale : la taille du lot}

Dans la perte NT-Xent (équation~\ref{eq:ntxent}), les négatifs sont les autres vues du lot.
À lot 256, chaque ancre est comparée à 510 négatifs. Les expériences de référence utilisent
des lots de plusieurs milliers d'images, soit un ordre de grandeur de plus, et
\cite{chen2020simclr} montre que la qualité des représentations croît nettement avec ce
nombre. Notre lot est dicté par la mémoire d'une carte unique de portable.

Cette limite n'a pas été mesurée : nous n'avons pas fait varier la taille du lot, faute de
budget. Nous la documentons parce qu'elle plafonne mécaniquement nos chiffres absolus et
interdit toute comparaison directe avec les résultats publiés. C'est aussi la piste
d'amélioration la plus rentable, et celle qui relie ce projet à la section~2 : MoCo
\cite{he2020moco} découple le nombre de négatifs de la taille du lot par une file d'attente,
BYOL \cite{grill2020byol} et DINO \cite{caron2021dino} suppriment le besoin de négatifs. À
matériel constant, l'accumulation de gradient ne résout pas le problème --- les négatifs
doivent être présents \emph{simultanément} dans le graphe de calcul.

\subsection{Une seule graine de pré-entraînement}

Toutes les incertitudes que nous rapportons proviennent des trois graines d'\emph{évaluation}.
Le pré-entraînement, lui, n'a tourné qu'une fois par configuration. La variabilité
pré-entraînement-à-pré-entraînement n'est donc pas quantifiée. Cela n'invalide pas les
écarts de plus de dix points, mais impose la prudence sur tout écart inférieur au point ---
d'où le traitement réservé au flip et au flou, et à l'écart de $+0{,}43$ ci-dessous.

\subsection{Le taux d'apprentissage a été réglé à 10\,\%, puis transféré}

Le balayage de la figure~\ref{fig:lr} a été conduit au régime $10\,\%$ des labels, et la
valeur retenue ($0{,}03$) a été appliquée aux trois régimes. Rien ne garantit que $0{,}03$
soit optimal à $100\,\%$ des labels pour la baseline supervisée ; une version antérieure de
cette baseline, réglée à $0{,}1$ sur l'ensemble du train, atteignait $94{,}79\,\%$. C'est la
raison pour laquelle nous parlons de \textbf{parité} à $100\,\%$ des labels et ne
revendiquons pas l'écart de $+0{,}43$ point comme une supériorité du pré-entraînement. Le
protocole reste équitable --- la même hypothèse de transfert pèse sur les deux méthodes ---
mais il est moins favorable à la baseline à ce régime précis, et il faut le dire.

\subsection{Le pré-entraînement a vu les images de validation}

Le pré-entraînement contrastif a consommé les 50\,000 images du train officiel, ce qui
inclut les 5\,000 images ensuite utilisées comme validation --- \emph{sans leurs labels}. Ce
n'est pas une fuite d'étiquettes, et c'est la pratique standard en SSL, où le
pré-entraînement est censé exploiter tout le corpus non annoté disponible. Nous le
mentionnons pour que le lecteur puisse en juger : la validation n'a servi qu'au choix du
taux d'apprentissage, et le test officiel n'a jamais été vu par le pré-entraînement.

\subsection{Autres limites}

\begin{itemize}
\item \textbf{Un seul jeu de données.} CIFAR-10 uniquement (section~\ref{sec:donnees}).
Rien ne garantit que les conclusions sur les augmentations se transfèrent à une autre
résolution ou à un autre domaine --- notre propre résultat sur le flou en est
l'illustration.
\item \textbf{Une seule architecture.} ResNet-18 adapté. L'effet de la profondeur ou de la
largeur de l'encodeur n'a pas été étudié.
\item \textbf{Ablation additive, non factorielle} (section~\ref{sec:ablation}).
\item \textbf{Pas de mesure de coût en temps} (section~\ref{sec:resultats}).
\end{itemize}

\subsection{Pistes d'amélioration}

Par ordre de rapport valeur/coût décroissant : (1) augmenter le nombre de négatifs, par une
file d'attente à la MoCo plutôt que par un lot plus grand, ce que le matériel autorise ;
(2) allonger le pré-entraînement à 400 ou 1\,000 époques, la courbe de la
figure~\ref{fig:preentrainement} n'étant pas plateau ; (3) répéter le pré-entraînement sur
trois graines pour quantifier l'incertitude qui manque aujourd'hui ; (4) reprendre le
protocole sur STL-10, dont la partition non étiquetée est conçue pour le SSL ; (5) comparer
à une méthode sans négatifs (BYOL) pour vérifier si l'avantage attendu se matérialise à
notre échelle.

% =====================================================================
\section{Répartition du travail}

Le tableau~\ref{tab:repartition} détaille les contributions de chaque membre. Les trois
rôles correspondent à ceux nommés par le sujet. Le principe retenu est que
\textbf{chaque membre possède le périmètre qu'il devra défendre à l'oral}. L'historique du
dépôt (\texttt{git shortlog -sn}) reflète cette répartition.

\begin{table}[htbp]
\centering\small
\caption{Répartition des rôles et des contributions.}
\label{tab:repartition}
\begin{tabularx}{\textwidth}{@{}llXc@{}}
\toprule
\textbf{Membre} & \textbf{Rôle (sujet)} & \textbf{Contributions précises} & \textbf{Sections} \\
\midrule
\membreKR{} & Data / Experiment Engineer
& Chargement et analyse exploratoire de CIFAR-10 ; chaîne d'augmentations et sa
paramétrisation par brique ; découpage stratifié train/validation et script de vérification
anti-fuite ; balayage symétrique du taux d'apprentissage ; documentation de reproduction
(\texttt{REPRODUCTION.md}), gestion des graines.
& 3, 4.2, 4.8, 4.9, 8 \\
\addlinespace
\membreAK{} & Model / Research Engineer
& Encodeur ResNet-18 adapté et tête de projection ; implémentation de la perte NT-Xent et
de sa vérification ; boucle de pré-entraînement (échauffement, cosinus, précision mixte) ;
les trois protocoles d'évaluation et le témoin aléatoire ; conception et exécution de
l'étude d'ablation.
& 4.1, 4.3--4.7, 5, 6 \\
\addlinespace
\membreAA{} & Reporting / Backend Developer
& Application web de démonstration (extraction des représentations, projections t-SNE et
PCA, plus proches voisins, matrice de confusion) ; README et documentation d'installation ;
figures ; assemblage du rapport et de la présentation.
& 1, 2, 7, 9 \\
\bottomrule
\end{tabularx}
\end{table}

% =====================================================================
\vspace{4pt}
\encadre{\textbf{Disponibilité du code et des données.} L'intégralité du code, les fichiers
de résultats et les instructions de reproduction sont publiés à l'adresse
\url{\lienDepot}. Chaque ligne des tableaux~\ref{tab:principal} et~\ref{tab:ablation} est
reproductible par une commande unique, listée en annexe~\ref{ann:commandes}.}

\begin{thebibliography}{9}
\small
\bibitem{chen2020simclr}
T.~Chen, S.~Kornblith, M.~Norouzi, G.~Hinton.
\emph{A Simple Framework for Contrastive Learning of Visual Representations}.
ICML, 2020. \url{https://arxiv.org/abs/2002.05709}

\bibitem{he2020moco}
K.~He, H.~Fan, Y.~Wu, S.~Xie, R.~Girshick.
\emph{Momentum Contrast for Unsupervised Visual Representation Learning}.
CVPR, 2020. \url{https://arxiv.org/abs/1911.05722}

\bibitem{grill2020byol}
J.-B.~Grill, F.~Strub, F.~Altché \emph{et al.}
\emph{Bootstrap Your Own Latent: A New Approach to Self-Supervised Learning}.
NeurIPS, 2020. \url{https://arxiv.org/abs/2006.07733}

\bibitem{caron2021dino}
M.~Caron, H.~Touvron, I.~Misra \emph{et al.}
\emph{Emerging Properties in Self-Supervised Vision Transformers}.
ICCV, 2021. \url{https://arxiv.org/abs/2104.14294}

\bibitem{he2016resnet}
K.~He, X.~Zhang, S.~Ren, J.~Sun.
\emph{Deep Residual Learning for Image Recognition}.
CVPR, 2016. \url{https://arxiv.org/abs/1512.03385}

\bibitem{krizhevsky2009cifar}
A.~Krizhevsky.
\emph{Learning Multiple Layers of Features from Tiny Images}.
Rapport technique, University of Toronto, 2009.
\url{https://www.cs.toronto.edu/~kriz/cifar.html}
\end{thebibliography}


% =====================================================================
%  ANNEXES — hors du quota de 8 à 20 pages
% =====================================================================
\appendix
\renewcommand{\thesection}{Annexe~\Alph{section}}
\titleformat{\section}{\large\bfseries\color{bleuP}}{\thesection{} ---}{0.5em}{}

\section{Hyperparamètres complets}
\label{ann:hyper}

Toutes les valeurs ci-dessous sont celles des exécutions publiées ; elles sont également
sérialisées dans les fichiers \texttt{resultats.json} du dépôt, ce qui permet de les
vérifier sans relire le code.

\begin{table}[htbp]
\centering\small
\caption{Pré-entraînement contrastif (\texttt{scripts/preentrainer\_experimental.py}).}
\begin{tabularx}{\textwidth}{@{}lXl@{}}
\toprule
\textbf{Paramètre} & \textbf{Valeur} & \textbf{Option} \\
\midrule
Données & 50\,000 images du train officiel, sans labels & \texttt{--data-dir} \\
Architecture & ResNet-18 adapté ($3\times3$ pas 1, sans max pooling) & --- \\
Représentation & $h \in \mathbb{R}^{512}$ & --- \\
Tête de projection & $512 \rightarrow 512 \rightarrow 128$, ReLU & --- \\
Époques & 100 & \texttt{--epochs} \\
Taille de lot & 256 (soit 512 vues, 510 négatifs par ancre) & \texttt{--batch-size} \\
Optimiseur & AdamW & --- \\
Taux d'apprentissage & $3 \cdot 10^{-4}$ & \texttt{--learning-rate} \\
Taux minimal & $10^{-6}$ & \texttt{--min-learning-rate} \\
Échauffement & 10 époques, linéaire & \texttt{--warmup-epochs} \\
Décroissance & cosinus sur les 90 époques restantes & --- \\
\emph{Weight decay} & $10^{-4}$ & \texttt{--weight-decay} \\
Température $\tau$ & $0{,}5$ & \texttt{--temperature} \\
Précision mixte & activée (\texttt{torch.amp}) & --- \\
Graine & 42 & \texttt{--seed} \\
\bottomrule
\end{tabularx}
\end{table}

\begin{table}[htbp]
\centering\small
\caption{Les trois protocoles d'évaluation aval. Le témoin à encodeur aléatoire suit
exactement le protocole de l'évaluation linéaire.}
\begin{tabularx}{\textwidth}{@{}lXXX@{}}
\toprule
& \textbf{Évaluation linéaire} & \textbf{Fine-tuning} & \textbf{Supervisé} \\
\midrule
Encodeur & gelé & entraîné & entraîné \\
Poids entraînables & $\approx 5\,130$ & $\approx 11{,}2$ M & $\approx 11{,}2$ M \\
Entrée du classifieur & $h$ standardisé & $h$ & $h$ \\
Époques & 50 & 100 & 100 \\
Taille de lot & 256 & 256 & 256 \\
Optimiseur & SGD, momentum $0{,}9$ & SGD, momentum $0{,}9$, Nesterov & SGD, momentum $0{,}9$, Nesterov \\
Taux d'apprentissage & $0{,}1$ & $0{,}03$ & $0{,}03$ \\
Décroissance & cosinus $\rightarrow \text{lr} \times 10^{-3}$ & cosinus $\rightarrow 10^{-4}$ & cosinus $\rightarrow 10^{-4}$ \\
\emph{Weight decay} & --- & $5 \cdot 10^{-4}$ & $5 \cdot 10^{-4}$ \\
Augmentations & aucune & \texttt{RandomCrop(32, pad=4)} $+$ \texttt{Flip} & \texttt{RandomCrop(32, pad=4)} $+$ \texttt{Flip} \\
Graines & 42, 123, 2026 & 42, 123, 2026 & 42, 123, 2026 \\
\bottomrule
\end{tabularx}
\end{table}

\section{Commandes de reproduction}
\label{ann:commandes}

Les commandes suivantes régénèrent chaque résultat publié. Elles supposent un environnement
Python conforme à \texttt{requirements.txt} et sont lancées depuis la racine du dépôt.

\begin{table}[htbp]
\centering\small
\caption{Correspondance résultat $\rightarrow$ commande $\rightarrow$ sortie.}
\begin{tabularx}{\textwidth}{@{}p{3.2cm}>{\raggedright\arraybackslash}X>{\raggedright\arraybackslash}p{3.3cm}@{}}
\toprule
\textbf{Résultat} & \textbf{Commande} & \textbf{Sortie} \\
\midrule
Pré-entraînement principal & \cmd{python -\brk{}m scripts.\brk{}preentrainer\_\brk{}experimental -\brk{}-\brk{}augmentations complet -\brk{}-\brk{}epochs 100} & \cmd{outputs/\brk{}simclr\_\brk{}cifar10\_\brk{}100ep/\brk{}} \\
\addlinespace
Tableau~\ref{tab:principal}, lignes \emph{linear probe} et \emph{témoin} & \cmd{python -\brk{}m scripts.\brk{}experimenter\_\brk{}fractions\_\brk{}labels -\brk{}-\brk{}seeds 42 123 2026 -\brk{}-\brk{}evaluer-\brk{}sur test} & \cmd{outputs/\brk{}experience\_\brk{}labels\_\brk{}45k/\brk{}} \\
\addlinespace
Tableau~\ref{tab:principal}, ligne \emph{fine-tuning} & \cmd{python -\brk{}m scripts.\brk{}experimenter\_\brk{}supervise\_\brk{}fractions -\brk{}-\brk{}checkpoint <ckpt> -\brk{}-\brk{}learning-\brk{}rate 0.\brk{}03 -\brk{}-\brk{}seeds 42 123 2026} & \cmd{outputs/\brk{}finetuning\_\brk{}simclr\_\brk{}45k/\brk{}} \\
\addlinespace
Tableau~\ref{tab:principal}, ligne \emph{supervisé} & \cmd{python -\brk{}m scripts.\brk{}experimenter\_\brk{}supervise\_\brk{}fractions -\brk{}-\brk{}learning-\brk{}rate 0.\brk{}03 -\brk{}-\brk{}seeds 42 123 2026} & \cmd{outputs/\brk{}supervise\_\brk{}from\_\brk{}scratch\_\brk{}45k/\brk{}} \\
\addlinespace
Tableau~\ref{tab:ablation} (intégral) & \cmd{python -\brk{}m scripts.\brk{}reproduire\_\brk{}ablation} & \cmd{outputs/\brk{}ablation/\brk{}tableau\_\brk{}ablation.\brk{}csv} \\
\addlinespace
Figure~\ref{fig:lr} (balayage) & \cmd{python -\brk{}m scripts.\brk{}experimenter\_\brk{}supervise\_\brk{}fractions -\brk{}-\brk{}learning-\brk{}rate <lr> -\brk{}-\brk{}evaluer-\brk{}sur validation} & \cmd{outputs/\brk{}*\_\brk{}recherche\_\brk{}lr/\brk{}} \\
\addlinespace
Contrôle anti-fuite & \cmd{python -\brk{}m scripts.\brk{}inspecter\_\brk{}split\_\brk{}validation} & \cmd{VERDICT : protocole propre} \\
\addlinespace
Application de démonstration & \cmd{python -\brk{}m scripts.\brk{}preparer\_\brk{}demo} & \cmd{demo/\brk{}index.\brk{}html} \\
\addlinespace
Figures du rapport & \cmd{python tools/\brk{}fig\_\brk{}resultats.\brk{}py} puis \cmd{tools/\brk{}fig\_\brk{}schemas.\brk{}py} & \cmd{docs/\brk{}rapport/\brk{}fig/\brk{}} \\
\bottomrule
\end{tabularx}
\end{table}

\section{Balayage du taux d'apprentissage, chiffres détaillés}
\label{ann:lr}

Accuracy sur les 5\,000 images de validation, au régime $10\,\%$ des labels, graine 42.
Ces chiffres sont ceux tracés en figure~\ref{fig:lr}.

\begin{table}[htbp]
\centering\small
\caption{Balayage symétrique. La même grille a été appliquée aux deux méthodes.}
\begin{tabular}{@{}lccccc@{}}
\toprule
\textbf{Taux d'apprentissage} & $0{,}003$ & $0{,}01$ & $\mathbf{0{,}03}$ & $0{,}1$ & $0{,}2$ \\
\midrule
Fine-tuning SimCLR      & $81{,}60$ & $81{,}44$ & $\mathbf{83{,}52}$ & $82{,}14$ & $63{,}46$ \\
Supervisé \emph{from scratch} & $60{,}60$ & $65{,}40$ & $\mathbf{70{,}46}$ & $68{,}84$ & $66{,}64$ \\
\midrule
Accuracy d'entraînement (fine-tuning) & $92{,}36$ & $99{,}73$ & $100{,}00$ & $100{,}00$ & $96{,}40$ \\
Accuracy d'entraînement (supervisé) & $84{,}42$ & $99{,}98$ & $99{,}98$ & $100{,}00$ & $94{,}42$ \\
\bottomrule
\end{tabular}
\end{table}

Les deux dernières lignes éclairent la forme des courbes. À $0{,}003$, aucune des deux
méthodes n'a fini de converger en 100 époques --- l'accuracy d'entraînement n'atteint pas
$100\,\%$. À $0{,}2$, elle redescend également : l'optimisation devient instable. Entre les
deux, les modèles atteignent la mémorisation complète et seule la généralisation les
départage. L'optimum commun à $0{,}03$ n'est donc pas un artefact de convergence.

\section{Coût de calcul}
\label{ann:cout}

\begin{table}[htbp]
\centering\small
\caption{Coût des cinq pré-entraînements, sur une NVIDIA GeForce RTX 4060 Laptop.
La colonne « durée totale observée » est donnée pour la transparence, non pour la
comparaison : voir la mise en garde ci-dessous.}
\begin{tabular}{@{}lcccc@{}}
\toprule
\textbf{Configuration} & \textbf{Époques} & \textbf{Médiane s/époque} & \textbf{Coût implicite}
& \textbf{Durée totale observée} \\
\midrule
crop seul               & 100 & $32{,}7$ & $\approx 55$ min & 914 min \\
crop $+$ flip           & 100 & $31{,}5$ & $\approx 53$ min & \phantom{0}53 min \\
$+$ color jitter        & 100 & $33{,}2$ & $\approx 55$ min & 457 min \\
$+$ niveaux de gris     & 100 & $30{,}4$ & $\approx 51$ min & \phantom{0}92 min \\
$+$ flou gaussien       & 100 & $29{,}7$ & $\approx 50$ min & 120 min \\
\bottomrule
\end{tabular}
\end{table}

Deux enseignements, et une mise en garde.

\paragraph{La chaîne d'augmentations ne coûte rien.} Les médianes s'échelonnent de $29{,}7$ à
$33{,}2$ secondes par époque, soit moins de $12\,\%$ d'écart entre la configuration la plus
légère et la plus lourde. Le coût du pré-entraînement est dominé par les passes avant et
arrière du ResNet, pas par les transformations d'images. Les cinq configurations de
l'ablation ont donc bien reçu, en pratique comme en principe, le même budget de calcul.

\paragraph{Le pic mémoire est identique partout} --- $1\,376$ Mo, à l'octet près pour les
cinq runs. C'est cohérent : la taille de lot et l'architecture sont fixes, et seules les
transformations, effectuées sur CPU, changent.

\paragraph{Mise en garde sur les durées totales.} Les durées totales observées vont de 53
minutes à 914 minutes pour des configurations dont le coût réel est identique à $12\,\%$
près. La machine était partagée avec d'autres charges pendant plusieurs runs, et quelques
époques isolées ont duré des dizaines de minutes. C'est la raison pour laquelle nous ne
publions \emph{aucune} comparaison de coût en temps entre méthodes
(section~\ref{sec:limites}) : la médiane par époque est la seule mesure de coût que notre
protocole soutient.

\paragraph{Budget total.} En s'appuyant sur les médianes, les cinq pré-entraînements
représentent environ $4{,}4$ heures de GPU. En y ajoutant les 36 exécutions d'évaluation
(4 bras $\times$ 3 fractions $\times$ 3 graines) et les 10 exécutions du balayage du taux
d'apprentissage, le budget total du projet se situe aux alentours de sept heures de calcul.

\section{Nomenclature}
\label{ann:nomenclature}

\begin{table}[htbp]
\centering\small
\begin{tabularx}{\textwidth}{@{}lX@{}}
\toprule
\textbf{Symbole ou terme} & \textbf{Signification} \\
\midrule
$x$ & une image du jeu de données, sans son étiquette \\
$\mathcal{T}$ & la famille d'augmentations aléatoires (recadrage, couleur, flou…) \\
$t, t'$ & deux transformations tirées indépendamment dans $\mathcal{T}$ \\
$\tilde{x}_i, \tilde{x}_j$ & les deux vues augmentées de $x$ ; elles forment une \emph{paire positive} \\
$f(\cdot)$ & l'encodeur — le ResNet-18 adapté, seul élément réutilisé en aval \\
$h \in \mathbb{R}^{512}$ & la représentation produite par l'encodeur \\
$g(\cdot)$ & la tête de projection — un MLP, jeté après le pré-entraînement \\
$z \in \mathbb{R}^{128}$ & la projection sur laquelle la perte contrastive s'applique \\
$\tau$ & la température de NT-Xent ; contrôle la dureté de la discrimination \\
$N$ & la taille de lot ; le lot fournit $2N-1$ candidats et $2N-2$ négatifs par ancre \\
\addlinespace
\emph{tâche prétexte} & la tâche artificielle résolue sans étiquettes (ici : reconnaître les deux vues d'une même image) \\
\emph{tâche aval} & la vraie tâche visée (ici : classer en 10 catégories) \\
\emph{évaluation linéaire} & protocole où l'encodeur est gelé et où seule une couche linéaire s'entraîne \\
\emph{fine-tuning} & protocole où l'encodeur et le classifieur s'adaptent ensemble \\
\emph{ablation} & retrait contrôlé d'un composant, tout le reste étant maintenu identique \\
\emph{témoin} & bras de contrôle servant à attribuer un effet à sa cause \\
\bottomrule
\end{tabularx}
\end{table}


\end{document}
